\textbf{Background.} Living kidney donors face a small but well-documented elevation in long-term risk of end-stage renal disease (ESRD), yet this risk has not been translated into a decision-relevant life-expectancy metric stratified by the demographic factors known to modify it.

\textbf{Methods.} I constructed a Markov microsimulation model comparing lifetime outcomes for matched donor and non-donor cohorts, entering at the age of donation and simulated in annual cycles, parameterized from the United States Renal Data System (USRDS), the Scientific Registry of Transplant Recipients (SRTR), and published cohort studies of post-donation ESRD incidence and mortality.

\textbf{Results.} In the base case, donation was associated with a mean life-expectancy cost of 911~days (2.5~years). This cost was slightly larger for Black donors than White donors and roughly 90\% larger at age~25 than at age~55. The result depends heavily on published estimates of the donor all-cause mortality hazard ratio (HR), which disagree not because they measure the same quantity with noise, but because the largest meta-analysis and the longest-follow-up cohort study describe different points on what is plausibly a non-proportional (time-varying) hazard, with excess mortality emerging only after approximately a decade. Holding the HR fixed at any single published value shifts the base-case cost from a net benefit of several years (HR~$=0.60$) to a cost of more than four years (the upper confidence bound of the long-follow-up estimate). Any other modeled parameter moves the result by a maximum of about six weeks. The life-expectancy value of prior-living-donor priority waitlist access and of a National-Kidney-Registry-style donation voucher was negligible unconditionally (under one day) but worth over one hundred days conditional on the designated beneficiary's actually developing ESRD.

\textbf{Conclusions.} These findings support age- and race-specific life-expectancy disclosure during living-donor informed consent, indicate that priority-access mechanisms function as catastrophic insurance rather than as interventions expected to shift average outcomes, and highlight resolving the time-course of donor all-cause mortality risk as a research priority, since it dominates every other source of uncertainty in this analysis.
